\chapter{माप और परिमाप}\label{ch:g4:measure}

पिछले साल हमने मापपट्टी, तराज़ू और घड़ी से मापा था
(\cref{ch:g3:measure}); यह अध्याय मात्रक-सारणी जोड़ता है ---
वह मशीन जो सब कुछ बदल देती है --- और चपटी आकृति नापने के दो
तरीक़े: उसकी सीमा की लंबाई (\omterm{def:g4:measure:perimeter}{परिमाप}) और उसके भीतर की जगह
(\omterm{def:g4:measure:area}{क्षेत्रफल}, फ़िलहाल खानों में गिना हुआ)।

\section{मात्रक-सारणी}

\begin{method}[सारणी से मात्रक बदलना]\label{met:g4:measure:table}
लंबाइयों के लिए हर स्तंभ अगले स्तंभ से दस गुना है:
\begin{center}
\renewcommand{\arraystretch}{1.3}
\begin{tabular}{c|c|c|c|c|c|c}
$\mathrm{km}$ & $\mathrm{hm}$ & $\mathrm{dam}$ & $\mathrm{m}$ & $\mathrm{dm}$ & $\mathrm{cm}$ & $\mathrm{mm}$ \\
\hline
 & & & $3$ & $4$ & $0$ & \\
\end{tabular}
\end{center}
\begin{enumerate}
  \item संख्या इस तरह लिखो कि हर स्तंभ में एक अंक आए, और माप का
        इकाई अंक अपने मात्रक के स्तंभ में हो (यहाँ $3.4$ m:
        $3$, m के नीचे);
  \item संख्या को नए मात्रक में पढ़ो: $3.4$ m $= 340$ cm
        (ख़ाली स्तंभ \omterm{def:g1:counting:zero}{शून्य} से भरो);
  \item \omterm{def:g2:measure:units}{द्रव्यमान} (kg\,\dots\,g) और \omterm{def:g2:measure:units}{धारिता} (L\,\dots\,cL) भी
        इसी तरह चलते हैं।
\end{enumerate}
\end{method}

\begin{example}\label{ex:g4:measure:convert}
$5$ km $= 5\,000$ m; \quad $270$ cm $= 2.7$ m ($2$, m के नीचे
आता है); \quad $3$ kg $250$ g $= 3\,250$ g; \quad
$2.5$ L $= 250$ cL।
\end{example}

\section{परिमाप}

\begin{definition}[परिमाप]\label{def:g4:measure:perimeter}
किसी \omterm{def:g4:geometry:polygon}{बहुभुज} का \emph{परिमाप}\index{परिमाप} उसकी सीमा की कुल
लंबाई है: उसकी सारी भुजाओं की लंबाइयाँ जोड़ दो (एक ही
मात्रक में!)। जानी-पहचानी आकृतियों के छोटे रास्ते:
\[
\text{आयत: } P = 2 \times (L + w),
\qquad
\text{वर्ग: } P = 4 \times c .
\]
\end{definition}

\begin{example}\label{ex:g4:measure:perimeter}
$12$ m गुणा $7$ m का आयताकार बग़ीचा:
\[
P = 2 \times (12 + 7) = 2 \times 19 = 38 \text{ m}
\]
बाड़ चाहिए। $18$ cm भुजा वाला वर्गाकार फ़ोटो-चौखटा:
$P = 4 \times 18 = 72$ cm।
\end{example}

\section{क्षेत्रफल: खाने गिनना}

\begin{definition}[गिनकर क्षेत्रफल]\label{def:g4:measure:area}
जाल पर बनी किसी आकृति का \emph{क्षेत्रफल}\index{क्षेत्रफल} उन
जाल-खानों की संख्या है जिन्हें वह ढकती है। दो आधे खाने मिलकर एक
गिने जाते हैं। क्षेत्रफल ``कितनी सतह'' का उत्तर देता है, \omterm{def:g4:measure:perimeter}{परिमाप}
``सीमा कितनी लंबी'' का --- दो अलग सवाल!
\end{definition}

\begin{omfigure}
\begin{tikzpicture}[scale=0.55]
  \draw[gray!40, thin] (0,0) grid (13,4);
  \draw[very thick, omDef, fill=omDef!15] (1,1) rectangle (5,3);
  \node[font=\small, omDef] at (3,2) {$8$ खाने};
  \draw[very thick, omThm, fill=omThm!15]
    (7,1) -- (11,1) -- (11,2) -- (9,2) -- (9,3) -- (7,3) -- cycle;
  \node[font=\small, omThm] at (8.4,1.5) {$6$ खाने};
\end{tikzpicture}

{\small जाल पर दो आकृतियाँ: \omterm{def:g4:measure:area}{क्षेत्रफल} $8$ और $6$ खाने। अब
सीमाओं पर चलकर इकाइयाँ गिनो: दोनों के लिए $12$! \omterm{def:g4:measure:perimeter}{परिमाप} वही,
\omterm{def:g4:measure:area}{क्षेत्रफल} अलग।}
\end{omfigure}

\begin{example}[आधे खाने]\label{ex:g4:measure:half}
जाल पर $4$ और $2$ भुजाओं वाला \omterm{def:g3:shapes:rightangle}{समकोण} त्रिभुज $3$ पूरे खाने और
$2$ आधे खाने ढकता है\,\dots\ ध्यान से गिनो: उसका \omterm{def:g4:measure:area}{क्षेत्रफल}
$4 \times 2$ के आयत का \omterm{def:g3:division:half}{आधा} है, यानी $4$ खाने। किसी आयत को उसके
विकर्ण पर काटने से हमेशा बराबर \omterm{def:g4:measure:area}{क्षेत्रफल} वाले दो त्रिभुज बनते
हैं।
\end{example}

\begin{example}[वही क्षेत्रफल, अलग परिमाप]\label{ex:g4:measure:samearea}
$4 \times 4$ का वर्ग और $8 \times 2$ का आयत दोनों $16$ खाने
ढकते हैं --- \omterm{def:g4:measure:area}{क्षेत्रफल} वही। सीमाएँ: वर्ग के लिए $16$ इकाई, आयत
के लिए $20$। कोई भी माप दूसरे को तय नहीं करता: लंबी-पतली आकृति
में थोड़ी-सी सतह के लिए बहुत सारी सीमा होती है।
\end{example}

\section{अवधियाँ}

\begin{method}[घंटे के पार समय जोड़ना]\label{met:g4:measure:time}
अवधियाँ जोड़नी या घटानी हों तो पूरे घंटों से होकर छलाँगों में
काम करो (याद रहे $1$ घंटा $= 60$ मिनट):
\[
9\text{ घं }40 + 35 \text{ मिनट}:
\quad 9\text{ घं }40 \xrightarrow{+20} 10\text{ घं }00
\xrightarrow{+15} 10\text{ घं }15 .
\]
मिनट $60$ के पार बिना बदले कभी मत जोड़ो: $40 + 35 = 75$ मिनट
$= 1$ घंटा $15$ मिनट।
\end{method}

\begin{example}\label{ex:g4:measure:timetable}
रेलगाड़ी $8{:}47$ पर चलती है और $11{:}05$ पर पहुँचती है। अवधि:
$8{:}47 \to 9{:}00$ पर $13$ मिनट; $9{:}00 \to 11{:}00$ पर
$2$ घंटे; $11{:}00 \to 11{:}05$ पर $5$ मिनट। कुल: $2$ घंटे $18$ मिनट।
\end{example}

\section{अभ्यास}

\begin{exercise}[$\star$]\label{exo:g4:measure:1}
मात्रक-सारणी से बदलो: $7$ m को cm में; $3.2$ km को m में;
$85$ mm को cm में; $4\,500$ m को km में।
\end{exercise}

\begin{exercise}[$\star$]\label{exo:g4:measure:2}
बदलो: $2$ kg को g में; $1\,800$ g को kg और g में; $3.5$ L को
cL में; $25$ cL को L में।
\end{exercise}

\begin{exercise}[$\star$]\label{exo:g4:measure:3}
\omterm{def:g4:measure:perimeter}{परिमाप} निकालो: $9$ cm $\times$ $5$ cm का आयत; $7.5$ cm भुजा
वाला वर्ग; $6$ cm, $8$ cm और $10$ cm भुजाओं वाला त्रिभुज।
\end{exercise}

\begin{exercise}[$\star$]\label{exo:g4:measure:4}
एक आयत का \omterm{def:g4:measure:perimeter}{परिमाप} $26$ cm और लंबाई $8$ cm है। उसकी चौड़ाई निकालो
(पहले लंबाई $+$ चौड़ाई निकालो)।
\end{exercise}

\begin{exercise}[$\star$]\label{exo:g4:measure:5}
जाल वाले काग़ज़ पर ठीक $10$ खाने ढकने वाली L-आकार की आकृति
बनाओ, और उसके \omterm{def:g4:measure:perimeter}{परिमाप} की इकाइयाँ गिनो।
\end{exercise}

\begin{exercise}[$\star$]\label{exo:g4:measure:6}
$12$ खानों के \omterm{def:g4:measure:area}{क्षेत्रफल} वाले दो अलग-अलग आयत बनाओ। हर एक का
\omterm{def:g4:measure:perimeter}{परिमाप} निकालो। कौन ``ज़्यादा गोल'' है, कौन ``ज़्यादा पतला''?
\end{exercise}

\begin{exercise}[$\star$]\label{exo:g4:measure:7}
जाल पर $6$ और $3$ भुजाओं वाला \omterm{def:g3:shapes:rightangle}{समकोण} त्रिभुज बना है। खानों में
उसका \omterm{def:g4:measure:area}{क्षेत्रफल} कितना है? (\cref{ex:g4:measure:half} का उपयोग करो।)
\end{exercise}

\begin{exercise}[$\star$]\label{exo:g4:measure:8}
अवधियाँ निकालो: $10{:}20$ से $12{:}45$ तक; $7{:}35$ से
$9{:}10$ तक। एक फ़िल्म $1$ घंटा $50$ मिनट की है और $20{:}30$ पर
शुरू होती है: वह कब ख़त्म होगी?
\end{exercise}

\begin{exercise}[$\star$]\label{exo:g4:measure:9}
किम $60$ m गुणा $45$ m के आयताकार मैदान के $3$ चक्कर दौड़ती है।
वह कितने मीटर दौड़ती है?
\end{exercise}

\begin{exercise}[$\star\star$]\label{exo:g4:measure:10}
एक किसान $85$ m भुजा वाले वर्गाकार खेत में बाड़ लगाना चाहता है,
और $4$ m का फाटक बिना बाड़ छोड़ना चाहता है। कितने मीटर बाड़ चाहिए?
\end{exercise}

\begin{exercise}[$\star\star$]\label{exo:g4:measure:11}
सही या ग़लत, जाल के उदाहरणों के साथ: ``दो आकृतियों का
\omterm{def:g4:measure:perimeter}{परिमाप} एक हो तो उनका \omterm{def:g4:measure:area}{क्षेत्रफल} भी एक होता है''; ``दो आकृतियों का
\omterm{def:g4:measure:area}{क्षेत्रफल} एक हो तो उनका \omterm{def:g4:measure:perimeter}{परिमाप} भी एक होता है''।
(\Cref{ex:g4:measure:samearea} और इस अध्याय की आकृति मदद करेगी।)
\end{exercise}
